% RESEARCH_DRAFT_SHA256: d43d78d1a74a235fa3c469677398fcc2ea219a30429d479f1175a4315e2747f8
% Layout and source-reference wrappers expanded for the web renderer.
\documentclass[11pt]{article}
\usepackage[margin=25mm]{geometry}
\usepackage{amsthm}
\usepackage{fontspec,unicode-math}
\setmainfont{Latin Modern Roman}
\setmathfont{Latin Modern Math}
\usepackage{amsmath,amssymb,mathtools}
\usepackage{xurl,hyperref}
\hypersetup{colorlinks=true,urlcolor=blue,linkcolor=blue}
\setcounter{secnumdepth}{0}
\setlength{\parindent}{0pt}
\setlength{\parskip}{5pt}
\setlength{\emergencystretch}{3em}
\newcommand{\sref}[2]{\hyperlink{source:#1:#2}{[S#2]}}
\newcommand{\cataloguescope}[1]{\paragraph{Source question.}#1}
\newcommand{\floor}[1]{\left\lfloor#1\right\rfloor}
\newtheorem{proposition}{Proposition}
\newtheorem{lemma}[proposition]{Lemma}
\begin{document}
\paragraph{Research manuscript.} OpenAI \texttt{gpt-6-astra}, reasoning \texttt{ultra}; 11 October 2026.
\paragraph{Preserved source body.} The complete catalogue mathematical text and sources follow with layout and source links expanded, including source qualifications and its historical preparation notes. New research begins at the section headed Research attempt.

\section{Semidefinite Hierarchies for Equiangular Line Bounds}\label{statement}

\subsection{Definitions and mathematical statement}

\paragraph{Shared notation.}Standard number systems are denoted by $\mathbb N,\mathbb Z,\mathbb Q,\mathbb R,\mathbb C$; index ranges, computational models and specialized notation follow the cited problem source and the definitions below.


For $0<a<1$, the graph $G_{n,a}$ has vertex set $S^{n-1}$; distinct $x,y$ are adjacent unless $x\cdot y\in\{a,-a\}$. Independent sets encode equiangular lines with common angle $\arccos(a)$. Write $I_k$ for independent sets of size at most $k$ and $I_{=j}$ for those of size $j$. A feasible kernel $T(x,y,Q)$, $Q\in I_{k-2}$, is symmetric in $x,y$ and has positive-semidefinite finite Gram matrices for every $Q$. Put $(B_kT)(S)=\sum_{Q\subseteq S,\,|Q|\leq k-2}\sum_{x,y\in S:\,Q\cup\{x,y\}=S}T(x,y,Q)$. The source's bound is $\Delta_k(G)^*=\inf(1+\lambda)$ subject to $B_kT\leq\lambda\mathbf1_{I_{=1}}-2\mathbf1_{I_{=2}}$ on nonempty $I_k$, using the continuous kernel space stated in the source. Here $\alpha(G)$ is the independence number.
\paragraph{Statement recorded in the supplied source.}
For the $k$-point semidefinite-programming bounds $\Delta_k(G_{n,a})^*$ for equiangular lines, explain their stabilization at $$(1/a^2-2)(1/a^2-1)/2,$$ determine whether $\Delta_k(G_{n,a})^*=\Theta(n^2)$ for fixed $k$ and $a$, and compare the corresponding level-$\alpha$ bound with the Lasserre hierarchy. \href{https://doi.org/10.4171/owr/2019/26}{[S2]}
\paragraph{Scope qualification.}
The parameter $a$ is the common inner-product magnitude; $\alpha$ denotes the graph independence number, not that magnitude. The stabilization target and comparison are those of the report's $k$-point hierarchy. The Lasserre hierarchy is the moment semidefinite hierarchy for the same independent-set problem; no equality between the hierarchies is asserted here.
\paragraph{Source status.}
The supplied catalogue labels this record open. Its source-status metadata was checked against the supplied snapshot on 2026-10-11. This import does not independently establish current open status or mathematical correctness. The catalogue literature assessment is dated 2026-08-22.
\subsection{Short English statement}
Explain observed stabilization and asymptotic growth of a semidefinite hierarchy bounding the number of equiangular lines.
\subsection{Sources}
\begin{itemize}
\item[S1] Ulam AI and the UnsolvedMath Contributors, UnsolvedMath: A Curated Collection of Open Mathematics Problems, catalogue entry. Catalogue curation and metadata: CC BY 4.0. The supplied source question is retained; definitions and the English summary are editorial additions.
\url{https://huggingface.co/datasets/ulamai/UnsolvedMath}.
\item[S2] Mixed-integer Nonlinear Optimization: a hatchery for modern mathematics (2020), Oberwolfach Reports, DOI 10.4171/owr/2019/26, pp. 1573--1637
\url{https://doi.org/10.4171/owr/2019/26}.
\item[S3] Mixed-integer Nonlinear Optimization, Oberwolfach Reports 26/2019, primary report PDF
\url{https://ems.press/content/serial-article-files/46805}.
\textit{Underlying Oberwolfach Reports material retains its source-specific licence, including CC BY-SA 4.0 where applicable; see THIRD\_PARTY\_NOTICES.md.}
\end{itemize}
\paragraph{Preparation and review.}
Original question attributed to the cited source and UnsolvedMath contributors. Codex (GPT-6; exact serving revision not exposed) assisted with definitions, English summaries, source inspection and formatting. No new human mathematical authorship or independent proof verification is asserted. Catalogue identity checked against the complete supplied registry and existing repository statements; known flagged repetitions and obvious identical conjectures excluded. Full mathematical-equivalence/literature review has not been performed. Statement preparation and final review followed contribution prompts C and D; automated checks verify structure and identity, not open status.
\label{end:statement}

\section{Research attempt}
\subsection{Literature, source context and scope}
The full original catalogue body and source qualifications are preserved above. The relevant primary workshop passages were read in the cached publisher report. Historical claims and subsequent attributed results are distinguished from the mathematical work below. A restricted example, standard lemma or known theorem is not counted as new progress on the entire original target. Missing assumptions and repeated or rhetorical extracts are not silently turned into different solved problems.

\subsection{The plateau number at one geometric endpoint}
The source asks three questions about its particular block-diagonal $k$-point hierarchy. An upper bound for actual line sets does not automatically give a feasible dual kernel for that hierarchy. We examine why its displayed number appears at a familiar geometric endpoint, and identify the missing inference.
Let $u_1,\ldots,u_N\in\mathbb R^n$ be unit representatives of distinct equiangular lines with $|u_i^Tu_j|=a$ for $i\ne j$, where $0<a<1$.
\begin{proposition}
One has $N\le n(n+1)/2$. If $na^2<1$, one also has
\[
N\le\frac{n(1-a^2)}{1-na^2}.
\]
At an integer $n_0=a^{-2}-2\ge1$, the second bound equals $n_0(n_0+1)/2=(a^{-2}-2)(a^{-2}-1)/2$.
\end{proposition}
\begin{proof}
The symmetric rank-one matrices $P_i=u_iu_i^T$ have Frobenius Gram matrix $(1-a^2)I_N+a^2J_N$, which is positive definite since its quadratic form is $(1-a^2)\|x\|^2+a^2(\sum x_i)^2>0$ for $x\ne0$. Hence the $P_i$ are linearly independent in the vector space of symmetric $n\times n$ matrices, whose dimension is $n(n+1)/2$.
For the second bound, the ordinary Gram matrix $G=(u_i^Tu_j)$ is PSD with rank at most $n$, trace $N$, and squared Frobenius norm $N+N(N-1)a^2$. Cauchy--Schwarz applied to its at most $n$ nonzero eigenvalues gives $N^2\le n[N+N(N-1)a^2]$. Divide by $N>0$ and rearrange, using $1-na^2>0$. The case $N=0$ is trivial. Substitution $a^2=1/(n_0+2)$ yields the last identity.
\end{proof}
For $a=1/3$, the endpoint $n_0=7$ is attained by 28 lines. For each two-element subset $\{i,j\}\subset\{1,\ldots,8\}$ set
\[
v_{ij}=e_i+e_j-\tfrac14\mathbf1_8.
\]
These vectors have coordinate sum zero and squared norm $3/2$. Distinct pair vectors have inner product $|\{i,j\}\cap\{k,\ell\}|-1/2$, hence magnitude $1/2$. After normalization the absolute inner product is $1/3$. Differences $v_{i\ell}-v_{j\ell}=e_i-e_j$, with $\ell$ distinct from $i,j$, span the seven-dimensional sum-zero subspace, so these lines do span dimension seven. Their number $\binom82=28$ equals both proved bounds.
\subsection{Verification and remaining gaps}
The 28 rational vectors' norms, absolute inner products and exact rank were checked. The failed inference would be to replace a bound on $M_a(n)$ by the identical bound on $\Delta_k(G_{n,a})^*$. The source only guarantees $M_a(n)\le\Delta_k(G_{n,a})^*$; an upper bound on the smaller quantity says nothing by itself about an upper bound on the larger one. No feasible kernel $T$ at the desired hierarchy level is constructed here. Thus the endpoint identity does not explain stabilization over other dimensions, the fixed-$k$ asymptotic behavior, or finite-level exactness. These classical geometric bounds and the known 28-line example are not claimed as new research results.

\section{References and source attribution}
Primary workshop report: \url{https://doi.org/10.4171/owr/2019/26}. Relevant passages read: 2019/26 pp. 1629--1632, block-diagonal hierarchy definition, numerical plateau table and three final questions.

All retrievals were checked on 11 October 2026. Original question attribution remains in the preserved Sources section. External theorems are attributed and not claimed as new discoveries or fully verified merely by citation.


\paragraph{Integration review.} Codex (GPT-6; exact serving revision not exposed) reviewed the written arguments and expanded layout and source-reference wrappers. Repository metadata and the production web renderer were checked using Python 3.13 and JavaScript with MathJax 3.2.2. This model-assisted review is not independent human proof verification; no PDF compilation or formal proof checking is asserted.

\section{Completion Estimate}
\noindent\textbf{COMPLETION: 0\%}
No substantiated new progress toward the entire original target is claimed. The worked results and failed route are delimited above; a known result, illustrative model, or auxiliary obstruction is not counted as original-target completion.

\section{Final status and attribution}
\textbf{Attempt status: unresolved.} This describes this research attempt, not a new ruling on the source catalogue's status or on external literature claims.
The mathematics and this self-audit were generated by OpenAI \texttt{gpt-6-astra}, reasoning \texttt{ultra}, on 11 October 2026. Codex (GPT-6; exact serving revision not exposed) orchestrated the batch, reviewed the written mathematical arguments, and contributed formatting/repository review. This additional model review is not independent human verification. No human mathematical authorship, independent human proof checking, or formal verification is claimed. The historical preparation paragraphs in the reproduced catalogue body concern that body only.

Contribution prompt A was applied to the complete definition and supplied source record, followed by the adversarial checks of prompt B. Local reproduction files are \texttt{batch022/checks.py}, \texttt{computations.json}, and the per-ID audit JSON. They were not public links at preparation time. Numerical checks are finite and do not prove unrestricted results. No LaTeX engine was found locally; compilation was not performed by this research agent.

\end{document}
